\chapter{Variabili aleatorie}

\Definizione{Variabile aleatoria}{
    Dato uno spazio di probabilità $(\Omega, \mathcal{F}, \mathbb{P})$, una \textbf{variabile aleatoria} su $\Omega$ è una funzione che ad ogni $\omega \in \Omega$ associa un numero reale
    \[
       \mathrm{ X: \Omega \rightarrow E \subset \mathbb{R}}
    \]
    tale che, per ogni $a \le b \in \mathbb{R}$, sia possibile calcolare la probabilità dell'insieme
    \[
        \mathrm{\{ \omega \in \Omega: a < X(\omega) \leq b \}}
    \]
    che, quindi, deve essere un evento di $\mathcal{F}$.
}

Le variabili aleatorie si distinguono in due categorie:
\begin{itemize}[label=$\triangleright$, itemsep=\smallskipamount]
    \item \textbf{variabili aleatorie discrete}
    \begin{itemize}[label=$\diamond$, itemsep=\smallskipamount]
        \item finite
        \item numerabili
    \end{itemize}
    \item \textbf{variabili aleatorie continue}
\end{itemize}

\Definizione{Variabili aleatorie discrete}{
    Una \textbf{variabile aleatoria} $X$ si dice \textbf{discreta} se l'insieme $E$ è discreto.
    
    In particolare:
    \begin{itemize}[label=$\circ$, itemsep=\smallskipamount]
        \item \textbf{finita}: se l'insieme $\mathrm{E}$ è un insieme finito
        \item \textbf{numerabile}: se l'insieme $\mathrm{E}$ è un insieme numerabile
    \end{itemize}
}

\Definizione{Variabile aleatoria continua}{
    Una \textbf{variabile aleatoria} si dice \textbf{continua} se l'insieme $\mathrm{E}$ è un insieme continuo.
}

\section{Variabili aleatorie finite}

\esempio{lancio del dado}{
    Consideriamo il lancio di un dado.
    Sia $\mathrm{X}$ la variabile che vale 0 se esce un numero pari e 1 se esce un numero dispari.
    
    \medskip
    Abbiamo così:
    \begin{itemize}[label=$\triangleright$, itemsep=\smallskipamount]
        \item $\Omega = \{1,2,3,4,5,6\}$
        \item $\mathrm{E = \{0,1\}}$ (insieme dei valori assunti da $\mathrm{X}$)
    \end{itemize}
    
    \smallskip
    Si ha:
    \[
        \begin{aligned}
            \mathrm{X(1) = X(3) = X(5)} &= \mathbf{1}\\
            \mathrm{X(2) = X(4) = X(6)} &= \mathbf{0}
        \end{aligned}
    \]
}

L'insieme dei valori assunti da $\mathrm{X}$ e le probabilità con cui vengono assunti si chiama \textbf{legge o distribuzione di $\boldsymbol{\mathrm{X}}$}.

\Definizione{Funzione di ripartizione}{
    La \textbf{funzione di ripartizione} di una variabile aleatoria discreta $\mathrm{X}$ è definita da:
    \[
        \boldsymbol{\mathrm{F_X(x) = \mathbb{P}(X \leq x) = \sum_{t \leq x} p_X(t)}}
    \]
}

La variabile alleatoria essendo una funzione di esperimenti, essa assume i suo valori con una certa probabilità.

\esempio{lancio del dado}{
    Consideriamo l'esempio fatto in precedenza sul lancio del dado.

    \smallskip
    \textcolor{red}{\textbf{Cosa vuol dire che $\boldsymbol{\mathrm{X = 0}}$?}}

    Significa che il lancio ha dato un risultato pari.
    \[
        \{\mathrm{X = 0}\} \equiv \{2,4,6\}
    \]

    Ciò significa che:
    \[
        \mathbb{P}(\mathrm{X = 0}) = \mathbb{P}(\{2,4,6\}) = \frac{3}{6} = \frac{1}{2}
    \]

    \medskip
    \textcolor{red}{\textbf{Cosa vuol dire che $\boldsymbol{\mathrm{X = 1}}$?}}

    Significa che il lancio ha dato un risultato dipari.
    \[
        \{\mathrm{X = 1}\} \equiv \{1,3,5\}
    \]

    Ciò significa che:
    \[
        \mathbb{P}(\mathrm{X = 1}) = \mathbb{P}(\{1,3,5\}) = \frac{3}{6} = \frac{1}{2}
    \]

    \medskip
    Facciamo una tabella riassuntiva:

    \begin{center}
        \renewcommand{\arraystretch}{1.5}
        \begin{tabular}{c|c|c}
            \textbf{Valore } $x_i$ & \textbf{Sottoinsieme di } $\Omega$ & $P(X = x_i)$ \\
            \hline
            $0$ & $\{2, 4, 6\}$ & $\frac{1}{2}$ \\
            \hline
            $1$ & $\{1, 3, 5\}$ & $\frac{1}{2}$ \\
        \end{tabular}
    \end{center}
}